\tituloI{PLANTEAMIENTO DEL ESTUDIO}

\tituloII{Planteamiento y formulación del problema}

\parrafoII{
Según \cite{smith2023global}, la gestión de inventarios en el sector minorista ha evolucionado hacia la integración de sistemas digitales que buscan garantizar la consistencia de la información y la trazabilidad de las operaciones. En este contexto, las organizaciones han adoptado soluciones tecnológicas orientadas a reducir la intervención manual en los procesos críticos. Sin embargo, a pesar de estos avances, la persistencia de errores humanos, junto con la ausencia de mecanismos de validación en tiempo real, continúa representando un desafío significativo. Estas deficiencias impactan directamente en la precisión de los registros de stock, generando discrepancias que derivan en pérdidas económicas, desabastecimiento de productos y deterioro en la calidad del servicio al cliente.
}

\parrafoII{
En el Perú, la digitalización de procesos comerciales ha sido impulsada por la obligatoriedad del uso de Comprobantes de Pago Electrónicos (CPE), lo que ha generado un crecimiento sostenido en la adopción de tecnologías digitales en las micro y pequeñas empresas. No obstante, como señala \cite{sunat2023informe}, esta adopción no ha sido homogénea ni completamente integrada, ya que muchas organizaciones continúan utilizando soluciones fragmentadas que dependen de proveedores externos para la facturación electrónica. Esta situación provoca la existencia de múltiples fuentes de datos que no se sincronizan adecuadamente, dificultando la gestión unificada de la información y aumentando el riesgo de inconsistencias operativas y tributarias.
}

\parrafoII{
A nivel local, la problemática se intensifica en los comercios de abarrotes de la región Cusco, donde las limitaciones de infraestructura tecnológica condicionan la implementación de soluciones robustas. De acuerdo con \cite{cusco2024comercio}, la digitalización en este sector se caracteriza por el uso de sistemas básicos o desarrollos informales que no responden a las exigencias de control y eficiencia operativa. En muchos casos, los registros de inventario se realizan de manera manual o mediante herramientas no integradas, lo que incrementa la probabilidad de errores y reduce la capacidad de respuesta ante la alta rotación de productos propia de este tipo de negocio.
}

\parrafoII{
Las causas de esta problemática pueden analizarse desde una perspectiva técnica. Tal como se expone en \cite{doe2021digital}, la presencia de silos de información en las pequeñas empresas limita la interoperabilidad entre sistemas y afecta la consistencia de los datos. En este sentido, la ausencia de restricciones de integridad a nivel de base de datos permite la ocurrencia de inconsistencias lógicas, como la generación de stocks negativos o registros incompletos. Asimismo, la dependencia de servicios externos de facturación introduce latencias y puntos de fallo adicionales, lo que debilita la arquitectura del sistema y reduce su confiabilidad.
}

\parrafoII{
Estas deficiencias generan múltiples efectos en la operación del negocio. Entre los principales impactos se encuentran la discrepancia entre el inventario físico y el lógico, la dificultad para planificar compras de manera eficiente y la pérdida de control sobre los flujos de mercancía. Además, la dependencia de sistemas externos para la emisión de comprobantes electrónicos provoca retrasos en el punto de venta, afectando la experiencia del cliente. Según \cite{garcia2023gestion}, estos problemas no solo repercuten en la eficiencia operativa, sino también en el cumplimiento de las obligaciones tributarias, incrementando el riesgo de sanciones y afectando la sostenibilidad del negocio.
}

\parrafoII{
Desde un enfoque prospectivo, la continuidad de estas condiciones podría generar una brecha competitiva cada vez mayor entre los pequeños comercios y las cadenas formalizadas que cuentan con sistemas integrados. La falta de control sobre la información y la dependencia tecnológica externa limitarían la capacidad de adaptación del negocio ante cambios en la normativa o en el mercado. En este sentido, resulta necesario evaluar soluciones basadas en arquitecturas locales con validaciones estrictas e integración directa de procesos críticos, que permitan mejorar la consistencia de los datos, optimizar el rendimiento del sistema y garantizar la autonomía operativa en entornos con recursos limitados.
}
\tituloII{Formulación del problema}

\tituloIII{Problema general}

\parrafoIII{¿De qué manera la implementación de un sistema de gestión local con integración de facturación electrónica directa e integridad de datos influye en la optimización del control operativo y la gestión comercial en un establecimiento de abarrotes en Cusco, 2026?}

\tituloIII{Problemas específicos}

\begin{listaPE}
\item ¿En qué medida la implementación de un motor de validación de existencias en tiempo real reduce la presencia de saldos negativos en el inventario?
\item ¿Cómo influye la integración de un módulo de facturación electrónica directa en la reducción de los tiempos de emisión de comprobantes en el punto de venta?
\item ¿De qué forma la centralización de datos en una arquitectura local mejora la confiabilidad de los reportes analíticos para la toma de decisiones?
\end{listaPE}

\tituloII{Objetivos}

\tituloIII{Objetivo general}

\parrafoIII{Evaluar el efecto de la implementación de un sistema de gestión local con facturación electrónica integrada y control de integridad de datos en la optimización de los procesos operativos y comerciales en un comercio de abarrotes en Cusco.}

\tituloIII{Objetivos específicos}

\begin{listaOE}
\item Evaluar la reducción de errores de inventario mediante la implementación de un motor de validación de transacciones.
\item Medir la disminución del tiempo de emisión de comprobantes electrónicos mediante la integración directa con los servicios web de SUNAT bajo el estándar UBL \cite{iso2024ubl}.
\item Analizar la mejora en la confiabilidad de los reportes de gestión a partir de la centralización de datos en una arquitectura local.
\end{listaOE}

\tituloII{Justificación de la investigación}

\tituloIII{Justificación teórica}

\parrafoIII{
La investigación se sustenta en principios de ingeniería de software relacionados con la integridad de datos, arquitectura de sistemas y gestión de transacciones. Aporta evidencia sobre la efectividad de las validaciones a nivel de base de datos como mecanismo para garantizar consistencia en sistemas de alta transaccionalidad. Asimismo, contribuye al análisis de arquitecturas locales frente a modelos desacoplados, especialmente en contextos donde los recursos tecnológicos son limitados \cite{garcia2023gestion}.
}

\tituloIII{Justificación práctica}

\parrafoIII{
La investigación aborda una problemática real en microempresas del sector retail, proporcionando una solución tecnológica que mejora el control de inventarios, reduce errores operativos y optimiza el proceso de facturación electrónica. Esto permite incrementar la eficiencia operativa, reducir costos asociados a errores y mejorar el cumplimiento tributario sin depender de servicios externos.
}

\tituloIII{Justificación metodológica}

\parrafoIII{
Se propone un enfoque metodológico basado en la implementación y evaluación de un sistema en un entorno real, permitiendo medir el impacto de variables técnicas sobre indicadores operativos. Además, se establece un marco replicable para el desarrollo de sistemas eficientes en hardware limitado bajo entornos Linux Bare Metal, contribuyendo a futuras investigaciones aplicadas \cite{garcia2023linux}.
}

\tituloII{Delimitación de la investigación}

\tituloIII{Delimitación espacial}

\parrafoIII{El estudio se realiza en un establecimiento minorista de abarrotes ubicado en la provincia y departamento de Cusco.}

\tituloIII{Delimitación temporal}

\parrafoIII{La investigación comprende el periodo de julio a septiembre del año 2026.}

\tituloIII{Delimitación social}

\parrafoIII{El estudio beneficia directamente al propietario y personal administrativo del establecimiento, al mejorar la gestión operativa y la confiabilidad de la información.}

\tituloII{Hipótesis y variables}

\tituloIII{Hipótesis}

\parrafoIII{La implementación de un sistema de gestión local con integridad de datos y facturación directa reduce los errores de inventario y disminuye el tiempo promedio de atención al cliente en el establecimiento de abarrotes.}

\tituloIII{Definición de variables}

\parrafoIII{Las variables del estudio son:}

\begin{listaVar}
\item \textbf{Variable independiente:} Implementación de un sistema de gestión local con validación de integridad y facturación electrónica directa.
\item \textbf{Variables dependientes:}
\begin{itemize}
\item Número de errores de inventario (stock negativo o inconsistente).
\item Tiempo promedio de emisión de comprobantes electrónicos.
\end{itemize}
\item \textbf{Variables intervinientes:} Capacidad del hardware local y estabilidad del entorno operativo.
\end{listaVar}